% ------------------------------------------------------------------
% FILE: sections/08_algorithm_review.tex
% Comparative review of 10 candidate algorithms for S11 classification.
% \input{} into main.tex — do NOT wrap in \documentclass or \begin{document}.
% ------------------------------------------------------------------

\section{Tiered Algorithm Strategy}\label{sec:algorithms}

Rather than conducting a blind flat search across arbitrary models, this methodology
adopts a principled \textbf{tiered complexity hierarchy} (Level 0 to Level 3). We
advance through complexity levels sequentially, only proceeding to the next tier if
the current tier fails to meet performance or robustness requirements.
Within this hierarchy, three architectures are highlighted as the primary candidates
for success: \textbf{Linear SVM / Logistic Regression} (for maximum interpretability
and minimal deployment cost), \textbf{Gradient Boosted Trees} (for consistent dominance
on tabular data and immunity to normalization tuning), and the \textbf{1D CNN} (for its
natural alignment with the local spectral structure of the data and excellent edge profile).

\paragraph{Review contract.}
To ensure that comparisons are fair and self-contained, every algorithm
review answers the same ten questions:
%
\begin{enumerate}
  \item \textbf{Core mechanism} --- what mathematical or structural principle underlies the model?
  \item \textbf{Expected input shape} --- what dimensionality and layout does the algorithm require?
  \item \textbf{Compatible representation route} --- Route~A (1D~spectral), Route~B (2D~image-like), or both?
  \item \textbf{Required preprocessing} --- what upstream pipeline stages are mandatory?
  \item \textbf{Preferred normalization} --- z-score, min--max, robust, or none?
  \item \textbf{Augmentation compatibility} --- which augmentation strategies from \Cref{sec:augment} apply?
  \item \textbf{Overfitting risk under \num{1603} samples} --- low, medium, or high?
  \item \textbf{Interpretability level} --- can predictions be traced back to physical frequency features?
  \item \textbf{Edge deployment suitability} --- model size, inference cost, and toolchain maturity.
  \item \textbf{Main failure modes} --- under what conditions will the algorithm break down?
\end{enumerate}



% ==================================================================
\subsection{Level 0: Unsupervised Structure Probe}\label{subsec:level0}
% ==================================================================
\subsubsection{K-Means Clustering}\label{subsec:kmeans}
% ==================================================================

\paragraph{Core mechanism.}
K-Means is an unsupervised centroid-based clustering algorithm that
partitions $n$ observations into $k$ clusters by iteratively minimising
the within-cluster sum of squared
distances~\cite{lloyd1982kmeans, macqueen1967kmeans}.  Given an
initial set of $k$ centroids, the algorithm alternates between
(1)~assigning each sample to the nearest centroid and (2)~updating
each centroid to the mean of its assigned samples, until convergence.

\paragraph{Critical framing.}
\textbf{K-Means is an unsupervised structure probe, not a primary
classifier.}  Its purpose in this pipeline is to answer a foundational
question \emph{before} supervised learning begins: do the three
physical states (\texttt{dry}, \texttt{wet}, \texttt{ice}) produce
naturally separable clusters in the chosen feature space?  If K-Means
with $k = 3$ recovers clusters that align well with the known labels,
this provides evidence that the feature representation captures
physically meaningful differences.  If not, the representation may
need revision before investing in supervised model training.

\paragraph{Expected input shape.}
Flattened 1D feature vectors, identical to the supervised methods.

\paragraph{Compatible route.}
Route~A.

\paragraph{Required preprocessing.}
Shared pipeline from \Cref{sec:preproc}.

\paragraph{Preferred normalization.}
\textbf{Z-score normalisation is essential.}  K-Means uses Euclidean
distance to assign cluster membership; without normalisation, features
with larger numerical ranges (e.g., phase in degrees vs.\ magnitude
in dB) will dominate the distance calculation and distort cluster
assignments.

\paragraph{Augmentation compatibility.}
Not applicable.  Augmentation is a supervised training technique;
K-Means operates on the original data to probe its natural structure.

\paragraph{Overfitting risk.}
\textbf{None in the conventional sense.}  K-Means does not use labels
during training, so there is no label leakage or label-driven
memorisation.  However, the algorithm \emph{can} find spurious
structure in high-dimensional spaces (the curse of dimensionality
facilitates apparent clustering even in random data).  Silhouette
analysis and comparison against random label permutations guard
against this.

\paragraph{Interpretability.}
\textbf{High.}  Each cluster centroid is a $d$-dimensional vector in
the same feature space as the input, and can be plotted directly as a
spectral trace.  This enables visual comparison: does the
\texttt{dry}-aligned centroid look like a prototypical dry-state
$S_{11}$ trace?  Cluster-to-class assignment can be evaluated via the
Adjusted Rand Index (ARI), which measures agreement between the
discovered partition and the true labels, corrected for chance.

\paragraph{Edge deployment suitability.}
\textbf{Very high.}  Inference is a nearest-centroid lookup: compute
$k = 3$ Euclidean distances and select the minimum.  This requires
$3 \times d$ multiply-accumulate operations and negligible memory
(three $d$-dimensional centroids).  The entire model fits in less than
\SI{100}{\kilo\byte} and can be implemented in any language.

\paragraph{Main failure modes.}
K-Means assumes that clusters are convex, isotropic, and roughly
equally sized --- assumptions that may not hold for $S_{11}$ feature
spaces.  If the \texttt{wet} and \texttt{ice} classes form elongated
or overlapping distributions, K-Means may (a)~split one true class
into two clusters, or (b)~merge two distinct classes into one cluster.
Running K-Means with multiple random initialisations, evaluating
silhouette scores, and comparing alternative values of $k$ (e.g.,
$k \in \{2, 3, 4, 5\}$) will help diagnose these failure modes.

\paragraph{Evaluation protocol.}
Because K-Means does not use labels, its output requires a distinct
evaluation protocol:
%
\begin{itemize}
  \item Run K-Means with $k = 3$ and multiple random initialisations
        (e.g., \num{50}).
  \item Assign each discovered cluster to the class label that has
        the highest overlap (Hungarian matching).
  \item Compute the Adjusted Rand Index (ARI) between the cluster
        assignments and the true labels.
  \item Compute the silhouette score to assess cluster cohesion and
        separation.
  \item Visualise cluster centroids as spectral traces and compare
        against per-class mean spectra.
  \item Repeat the analysis for each candidate feature encoding
        (magnitude-only, magnitude-plus-phase, etc.) to determine
        which representation yields the cleanest unsupervised
        separation.
\end{itemize}


% ==================================================================
% Summary comparison table (defined in tables/algorithm_comparison.tex)
% ==================================================================

\bigskip
% tables/algorithm_comparison.tex — Master algorithm comparison table
% Included via % tables/algorithm_comparison.tex — Master algorithm comparison table
% Included via % tables/algorithm_comparison.tex — Master algorithm comparison table
% Included via \input{tables/algorithm_comparison} from a section file.
% Designed for landscape orientation; wrap the \input{} call inside
% a \begin{landscape}...\end{landscape} environment (from pdflscape) or
% place it on a standalone landscape page.

\begin{footnotesize}
\setlength{\LTcapwidth}{\linewidth}
\setlength{\tabcolsep}{3.5pt}
\renewcommand{\arraystretch}{1.25}

\begin{longtable}{%
  >{\raggedright\arraybackslash}p{2.1cm}   % Algorithm
  >{\raggedright\arraybackslash}p{1.5cm}   % Family
  >{\raggedright\arraybackslash}p{1.8cm}   % Input Shape
  >{\centering\arraybackslash}p{1.0cm}     % Route
  >{\raggedright\arraybackslash}p{1.6cm}   % Normalization
  >{\raggedright\arraybackslash}p{1.7cm}   % Augmentation
  >{\centering\arraybackslash}p{1.1cm}     % Overfit Risk
  >{\centering\arraybackslash}p{1.1cm}     % Interpretability
  >{\centering\arraybackslash}p{1.0cm}     % Edge Suit.
  >{\centering\arraybackslash}p{1.0cm}     % Noise Rob.
  >{\raggedright\arraybackslash}p{2.8cm}   % Key Failure Mode
}
\caption{Qualitative comparison of candidate algorithms for three-class $S_{11}$
         classification (\texttt{dry}/\texttt{wet}/\texttt{ice}).
         Assessments reflect expected behaviour given the dataset size
         ($N=1603$) and are \emph{not} derived from experimental results.}
\label{tab:algorithm_comparison} \\

\toprule
\textbf{Algorithm} &
\textbf{Family} &
\textbf{Input Shape} &
\textbf{Route} &
\textbf{Preferred Norm.} &
\textbf{Augment.\ Compat.} &
\textbf{Overfit Risk} &
\textbf{Interp.} &
\textbf{Edge Suit.} &
\textbf{Noise Rob.} &
\textbf{Key Failure Mode} \\
\midrule
\endfirsthead

% ── continuation header ──
\multicolumn{11}{l}{\small\itshape (continued from previous page)} \\[2pt]
\toprule
\textbf{Algorithm} &
\textbf{Family} &
\textbf{Input Shape} &
\textbf{Route} &
\textbf{Preferred Norm.} &
\textbf{Augment.\ Compat.} &
\textbf{Overfit Risk} &
\textbf{Interp.} &
\textbf{Edge Suit.} &
\textbf{Noise Rob.} &
\textbf{Key Failure Mode} \\
\midrule
\endhead

% ── footer ──
\midrule
\multicolumn{11}{r}{\small\itshape (continued on next page)} \\
\endfoot

% ── last footer ──
\bottomrule
\endlastfoot

% ═══════════════════════════════════════════════════════════════
%  ROW 1 — Logistic Regression
% ═══════════════════════════════════════════════════════════════
Logistic Regression
  & Linear
  & $(D,)$ flat vector
  & A
  & StandardScaler (zero-mean, unit-var)
  & Limited; noise injection only
  & Low
  & High
  & High
  & Low
  & Under-fits if class boundaries are non-linear in spectral space \\

% ═══════════════════════════════════════════════════════════════
%  ROW 2 — Linear SVM
% ═══════════════════════════════════════════════════════════════
Linear SVM
  & Linear
  & $(D,)$ flat vector
  & A
  & StandardScaler
  & Limited; noise injection only
  & Low
  & High
  & High
  & Low--Med
  & Same linear limitation; sensitive to outlier support vectors \\

% ═══════════════════════════════════════════════════════════════
%  ROW 3 — RBF SVM
% ═══════════════════════════════════════════════════════════════
RBF SVM
  & Kernel
  & $(D,)$ flat vector
  & A
  & StandardScaler (critical)
  & Limited; noise injection, SMOTE
  & Medium
  & Low
  & Medium
  & Medium
  & $\gamma$--$C$ grid search costly; kernel matrix scales $O(N^2)$ \\

% ═══════════════════════════════════════════════════════════════
%  ROW 4 — Random Forest
% ═══════════════════════════════════════════════════════════════
Random Forest
  & Ensemble (bagging)
  & $(D,)$ flat vector
  & A
  & Not required (tree-based)
  & SMOTE, noise injection
  & Low--Med
  & Medium
  & High
  & High
  & Can memorise noisy features if trees are too deep; feature importance may
    mislead with correlated spectral bins \\

% ═══════════════════════════════════════════════════════════════
%  ROW 5 — Gradient Boosted Trees (XGBoost)
% ═══════════════════════════════════════════════════════════════
Gradient Boosted Trees (XGBoost)
  & Ensemble (boosting)
  & $(D,)$ flat vector
  & A
  & Not required (tree-based)
  & SMOTE, noise injection
  & Medium
  & Medium
  & High
  & High
  & Prone to over-fitting with many boosting rounds on small $N$;
    requires careful early stopping \\

% ═══════════════════════════════════════════════════════════════
%  ROW 6 — Multilayer Perceptron
% ═══════════════════════════════════════════════════════════════
Multilayer Perceptron
  & Neural network (feedforward)
  & $(D,)$ flat vector
  & A
  & StandardScaler or MinMax
  & Noise injection, mixup, dropout
  & High
  & Low
  & Medium
  & Medium
  & Over-parameterised for small $N$; sensitive to learning rate and
    architecture choices \\

% ═══════════════════════════════════════════════════════════════
%  ROW 7 — 1D CNN
% ═══════════════════════════════════════════════════════════════
1D CNN
  & Neural network (convolutional)
  & $(L, C)$ sequence
  & A
  & Per-channel StandardScaler or BatchNorm
  & Noise injection, frequency masking, mixup
  & High
  & Low
  & Medium
  & Med--High
  & Needs sufficient depth for receptive field; over-fits quickly on 1603
    samples without strong regularisation \\

% ═══════════════════════════════════════════════════════════════
%  ROW 8 — ResNet-1D
% ═══════════════════════════════════════════════════════════════
ResNet-1D
  & Neural network (residual)
  & $(L, C)$ sequence
  & A
  & BatchNorm (built-in)
  & Noise injection, frequency masking, mixup
  & High
  & Low
  & Low--Med
  & Med--High
  & Highest parameter count; extremely prone to over-fitting at $N=1603$;
    large model footprint challenges edge deployment \\

% ═══════════════════════════════════════════════════════════════
%  ROW 9 — KAN
% ═══════════════════════════════════════════════════════════════
Kolmogorov--Arnold Network (KAN)
  & Neural network (spline-based)
  & $(D,)$ flat vector
  & A
  & StandardScaler
  & Noise injection, mixup
  & Med--High
  & Medium
  & Medium
  & Medium
  & Immature tooling; spline knot count acts as hidden capacity control;
    limited edge-runtime support \\

% ═══════════════════════════════════════════════════════════════
%  ROW 10 — K-Means (unsupervised)
% ═══════════════════════════════════════════════════════════════
K-Means\textsuperscript{\dag}
  & Unsupervised (centroid)
  & $(D,)$ flat vector
  & A
  & MinMaxScaler or StandardScaler
  & N/A (unsupervised)
  & Low
  & Medium
  & High
  & Low
  & Assumes spherical clusters; label assignment is post-hoc; cannot capture
    non-convex class boundaries \\

\end{longtable}

\vspace{0.3em}
\noindent\textsuperscript{\dag}\,K-Means is included as an \emph{unsupervised probe}
to assess natural cluster separability in the feature space, not as a
direct classification algorithm.  Cluster-to-class mapping is evaluated
via majority voting and the adjusted Rand index.

\end{footnotesize}
 from a section file.
% Designed for landscape orientation; wrap the \input{} call inside
% a \begin{landscape}...\end{landscape} environment (from pdflscape) or
% place it on a standalone landscape page.

\begin{footnotesize}
\setlength{\LTcapwidth}{\linewidth}
\setlength{\tabcolsep}{3.5pt}
\renewcommand{\arraystretch}{1.25}

\begin{longtable}{%
  >{\raggedright\arraybackslash}p{2.1cm}   % Algorithm
  >{\raggedright\arraybackslash}p{1.5cm}   % Family
  >{\raggedright\arraybackslash}p{1.8cm}   % Input Shape
  >{\centering\arraybackslash}p{1.0cm}     % Route
  >{\raggedright\arraybackslash}p{1.6cm}   % Normalization
  >{\raggedright\arraybackslash}p{1.7cm}   % Augmentation
  >{\centering\arraybackslash}p{1.1cm}     % Overfit Risk
  >{\centering\arraybackslash}p{1.1cm}     % Interpretability
  >{\centering\arraybackslash}p{1.0cm}     % Edge Suit.
  >{\centering\arraybackslash}p{1.0cm}     % Noise Rob.
  >{\raggedright\arraybackslash}p{2.8cm}   % Key Failure Mode
}
\caption{Qualitative comparison of candidate algorithms for three-class $S_{11}$
         classification (\texttt{dry}/\texttt{wet}/\texttt{ice}).
         Assessments reflect expected behaviour given the dataset size
         ($N=1603$) and are \emph{not} derived from experimental results.}
\label{tab:algorithm_comparison} \\

\toprule
\textbf{Algorithm} &
\textbf{Family} &
\textbf{Input Shape} &
\textbf{Route} &
\textbf{Preferred Norm.} &
\textbf{Augment.\ Compat.} &
\textbf{Overfit Risk} &
\textbf{Interp.} &
\textbf{Edge Suit.} &
\textbf{Noise Rob.} &
\textbf{Key Failure Mode} \\
\midrule
\endfirsthead

% ── continuation header ──
\multicolumn{11}{l}{\small\itshape (continued from previous page)} \\[2pt]
\toprule
\textbf{Algorithm} &
\textbf{Family} &
\textbf{Input Shape} &
\textbf{Route} &
\textbf{Preferred Norm.} &
\textbf{Augment.\ Compat.} &
\textbf{Overfit Risk} &
\textbf{Interp.} &
\textbf{Edge Suit.} &
\textbf{Noise Rob.} &
\textbf{Key Failure Mode} \\
\midrule
\endhead

% ── footer ──
\midrule
\multicolumn{11}{r}{\small\itshape (continued on next page)} \\
\endfoot

% ── last footer ──
\bottomrule
\endlastfoot

% ═══════════════════════════════════════════════════════════════
%  ROW 1 — Logistic Regression
% ═══════════════════════════════════════════════════════════════
Logistic Regression
  & Linear
  & $(D,)$ flat vector
  & A
  & StandardScaler (zero-mean, unit-var)
  & Limited; noise injection only
  & Low
  & High
  & High
  & Low
  & Under-fits if class boundaries are non-linear in spectral space \\

% ═══════════════════════════════════════════════════════════════
%  ROW 2 — Linear SVM
% ═══════════════════════════════════════════════════════════════
Linear SVM
  & Linear
  & $(D,)$ flat vector
  & A
  & StandardScaler
  & Limited; noise injection only
  & Low
  & High
  & High
  & Low--Med
  & Same linear limitation; sensitive to outlier support vectors \\

% ═══════════════════════════════════════════════════════════════
%  ROW 3 — RBF SVM
% ═══════════════════════════════════════════════════════════════
RBF SVM
  & Kernel
  & $(D,)$ flat vector
  & A
  & StandardScaler (critical)
  & Limited; noise injection, SMOTE
  & Medium
  & Low
  & Medium
  & Medium
  & $\gamma$--$C$ grid search costly; kernel matrix scales $O(N^2)$ \\

% ═══════════════════════════════════════════════════════════════
%  ROW 4 — Random Forest
% ═══════════════════════════════════════════════════════════════
Random Forest
  & Ensemble (bagging)
  & $(D,)$ flat vector
  & A
  & Not required (tree-based)
  & SMOTE, noise injection
  & Low--Med
  & Medium
  & High
  & High
  & Can memorise noisy features if trees are too deep; feature importance may
    mislead with correlated spectral bins \\

% ═══════════════════════════════════════════════════════════════
%  ROW 5 — Gradient Boosted Trees (XGBoost)
% ═══════════════════════════════════════════════════════════════
Gradient Boosted Trees (XGBoost)
  & Ensemble (boosting)
  & $(D,)$ flat vector
  & A
  & Not required (tree-based)
  & SMOTE, noise injection
  & Medium
  & Medium
  & High
  & High
  & Prone to over-fitting with many boosting rounds on small $N$;
    requires careful early stopping \\

% ═══════════════════════════════════════════════════════════════
%  ROW 6 — Multilayer Perceptron
% ═══════════════════════════════════════════════════════════════
Multilayer Perceptron
  & Neural network (feedforward)
  & $(D,)$ flat vector
  & A
  & StandardScaler or MinMax
  & Noise injection, mixup, dropout
  & High
  & Low
  & Medium
  & Medium
  & Over-parameterised for small $N$; sensitive to learning rate and
    architecture choices \\

% ═══════════════════════════════════════════════════════════════
%  ROW 7 — 1D CNN
% ═══════════════════════════════════════════════════════════════
1D CNN
  & Neural network (convolutional)
  & $(L, C)$ sequence
  & A
  & Per-channel StandardScaler or BatchNorm
  & Noise injection, frequency masking, mixup
  & High
  & Low
  & Medium
  & Med--High
  & Needs sufficient depth for receptive field; over-fits quickly on 1603
    samples without strong regularisation \\

% ═══════════════════════════════════════════════════════════════
%  ROW 8 — ResNet-1D
% ═══════════════════════════════════════════════════════════════
ResNet-1D
  & Neural network (residual)
  & $(L, C)$ sequence
  & A
  & BatchNorm (built-in)
  & Noise injection, frequency masking, mixup
  & High
  & Low
  & Low--Med
  & Med--High
  & Highest parameter count; extremely prone to over-fitting at $N=1603$;
    large model footprint challenges edge deployment \\

% ═══════════════════════════════════════════════════════════════
%  ROW 9 — KAN
% ═══════════════════════════════════════════════════════════════
Kolmogorov--Arnold Network (KAN)
  & Neural network (spline-based)
  & $(D,)$ flat vector
  & A
  & StandardScaler
  & Noise injection, mixup
  & Med--High
  & Medium
  & Medium
  & Medium
  & Immature tooling; spline knot count acts as hidden capacity control;
    limited edge-runtime support \\

% ═══════════════════════════════════════════════════════════════
%  ROW 10 — K-Means (unsupervised)
% ═══════════════════════════════════════════════════════════════
K-Means\textsuperscript{\dag}
  & Unsupervised (centroid)
  & $(D,)$ flat vector
  & A
  & MinMaxScaler or StandardScaler
  & N/A (unsupervised)
  & Low
  & Medium
  & High
  & Low
  & Assumes spherical clusters; label assignment is post-hoc; cannot capture
    non-convex class boundaries \\

\end{longtable}

\vspace{0.3em}
\noindent\textsuperscript{\dag}\,K-Means is included as an \emph{unsupervised probe}
to assess natural cluster separability in the feature space, not as a
direct classification algorithm.  Cluster-to-class mapping is evaluated
via majority voting and the adjusted Rand index.

\end{footnotesize}
 from a section file.
% Designed for landscape orientation; wrap the \input{} call inside
% a \begin{landscape}...\end{landscape} environment (from pdflscape) or
% place it on a standalone landscape page.

\begin{footnotesize}
\setlength{\LTcapwidth}{\linewidth}
\setlength{\tabcolsep}{3.5pt}
\renewcommand{\arraystretch}{1.25}

\begin{longtable}{%
  >{\raggedright\arraybackslash}p{2.1cm}   % Algorithm
  >{\raggedright\arraybackslash}p{1.5cm}   % Family
  >{\raggedright\arraybackslash}p{1.8cm}   % Input Shape
  >{\centering\arraybackslash}p{1.0cm}     % Route
  >{\raggedright\arraybackslash}p{1.6cm}   % Normalization
  >{\raggedright\arraybackslash}p{1.7cm}   % Augmentation
  >{\centering\arraybackslash}p{1.1cm}     % Overfit Risk
  >{\centering\arraybackslash}p{1.1cm}     % Interpretability
  >{\centering\arraybackslash}p{1.0cm}     % Edge Suit.
  >{\centering\arraybackslash}p{1.0cm}     % Noise Rob.
  >{\raggedright\arraybackslash}p{2.8cm}   % Key Failure Mode
}
\caption{Qualitative comparison of candidate algorithms for three-class $S_{11}$
         classification (\texttt{dry}/\texttt{wet}/\texttt{ice}).
         Assessments reflect expected behaviour given the dataset size
         ($N=1603$) and are \emph{not} derived from experimental results.}
\label{tab:algorithm_comparison} \\

\toprule
\textbf{Algorithm} &
\textbf{Family} &
\textbf{Input Shape} &
\textbf{Route} &
\textbf{Preferred Norm.} &
\textbf{Augment.\ Compat.} &
\textbf{Overfit Risk} &
\textbf{Interp.} &
\textbf{Edge Suit.} &
\textbf{Noise Rob.} &
\textbf{Key Failure Mode} \\
\midrule
\endfirsthead

% ── continuation header ──
\multicolumn{11}{l}{\small\itshape (continued from previous page)} \\[2pt]
\toprule
\textbf{Algorithm} &
\textbf{Family} &
\textbf{Input Shape} &
\textbf{Route} &
\textbf{Preferred Norm.} &
\textbf{Augment.\ Compat.} &
\textbf{Overfit Risk} &
\textbf{Interp.} &
\textbf{Edge Suit.} &
\textbf{Noise Rob.} &
\textbf{Key Failure Mode} \\
\midrule
\endhead

% ── footer ──
\midrule
\multicolumn{11}{r}{\small\itshape (continued on next page)} \\
\endfoot

% ── last footer ──
\bottomrule
\endlastfoot

% ═══════════════════════════════════════════════════════════════
%  ROW 1 — Logistic Regression
% ═══════════════════════════════════════════════════════════════
Logistic Regression
  & Linear
  & $(D,)$ flat vector
  & A
  & StandardScaler (zero-mean, unit-var)
  & Limited; noise injection only
  & Low
  & High
  & High
  & Low
  & Under-fits if class boundaries are non-linear in spectral space \\

% ═══════════════════════════════════════════════════════════════
%  ROW 2 — Linear SVM
% ═══════════════════════════════════════════════════════════════
Linear SVM
  & Linear
  & $(D,)$ flat vector
  & A
  & StandardScaler
  & Limited; noise injection only
  & Low
  & High
  & High
  & Low--Med
  & Same linear limitation; sensitive to outlier support vectors \\

% ═══════════════════════════════════════════════════════════════
%  ROW 3 — RBF SVM
% ═══════════════════════════════════════════════════════════════
RBF SVM
  & Kernel
  & $(D,)$ flat vector
  & A
  & StandardScaler (critical)
  & Limited; noise injection, SMOTE
  & Medium
  & Low
  & Medium
  & Medium
  & $\gamma$--$C$ grid search costly; kernel matrix scales $O(N^2)$ \\

% ═══════════════════════════════════════════════════════════════
%  ROW 4 — Random Forest
% ═══════════════════════════════════════════════════════════════
Random Forest
  & Ensemble (bagging)
  & $(D,)$ flat vector
  & A
  & Not required (tree-based)
  & SMOTE, noise injection
  & Low--Med
  & Medium
  & High
  & High
  & Can memorise noisy features if trees are too deep; feature importance may
    mislead with correlated spectral bins \\

% ═══════════════════════════════════════════════════════════════
%  ROW 5 — Gradient Boosted Trees (XGBoost)
% ═══════════════════════════════════════════════════════════════
Gradient Boosted Trees (XGBoost)
  & Ensemble (boosting)
  & $(D,)$ flat vector
  & A
  & Not required (tree-based)
  & SMOTE, noise injection
  & Medium
  & Medium
  & High
  & High
  & Prone to over-fitting with many boosting rounds on small $N$;
    requires careful early stopping \\

% ═══════════════════════════════════════════════════════════════
%  ROW 6 — Multilayer Perceptron
% ═══════════════════════════════════════════════════════════════
Multilayer Perceptron
  & Neural network (feedforward)
  & $(D,)$ flat vector
  & A
  & StandardScaler or MinMax
  & Noise injection, mixup, dropout
  & High
  & Low
  & Medium
  & Medium
  & Over-parameterised for small $N$; sensitive to learning rate and
    architecture choices \\

% ═══════════════════════════════════════════════════════════════
%  ROW 7 — 1D CNN
% ═══════════════════════════════════════════════════════════════
1D CNN
  & Neural network (convolutional)
  & $(L, C)$ sequence
  & A
  & Per-channel StandardScaler or BatchNorm
  & Noise injection, frequency masking, mixup
  & High
  & Low
  & Medium
  & Med--High
  & Needs sufficient depth for receptive field; over-fits quickly on 1603
    samples without strong regularisation \\

% ═══════════════════════════════════════════════════════════════
%  ROW 8 — ResNet-1D
% ═══════════════════════════════════════════════════════════════
ResNet-1D
  & Neural network (residual)
  & $(L, C)$ sequence
  & A
  & BatchNorm (built-in)
  & Noise injection, frequency masking, mixup
  & High
  & Low
  & Low--Med
  & Med--High
  & Highest parameter count; extremely prone to over-fitting at $N=1603$;
    large model footprint challenges edge deployment \\

% ═══════════════════════════════════════════════════════════════
%  ROW 9 — KAN
% ═══════════════════════════════════════════════════════════════
Kolmogorov--Arnold Network (KAN)
  & Neural network (spline-based)
  & $(D,)$ flat vector
  & A
  & StandardScaler
  & Noise injection, mixup
  & Med--High
  & Medium
  & Medium
  & Medium
  & Immature tooling; spline knot count acts as hidden capacity control;
    limited edge-runtime support \\

% ═══════════════════════════════════════════════════════════════
%  ROW 10 — K-Means (unsupervised)
% ═══════════════════════════════════════════════════════════════
K-Means\textsuperscript{\dag}
  & Unsupervised (centroid)
  & $(D,)$ flat vector
  & A
  & MinMaxScaler or StandardScaler
  & N/A (unsupervised)
  & Low
  & Medium
  & High
  & Low
  & Assumes spherical clusters; label assignment is post-hoc; cannot capture
    non-convex class boundaries \\

\end{longtable}

\vspace{0.3em}
\noindent\textsuperscript{\dag}\,K-Means is included as an \emph{unsupervised probe}
to assess natural cluster separability in the feature space, not as a
direct classification algorithm.  Cluster-to-class mapping is evaluated
via majority voting and the adjusted Rand index.

\end{footnotesize}


% ==================================================================
% Framework clarification
% ==================================================================

\paragraph{A note on implementation frameworks.}
Keras, TensorFlow, and PyTorch are \emph{software frameworks} for
implementing and training neural network architectures --- they are
not algorithms in themselves.  Keras, for example, is a high-level API
that can be used to implement the MLP (\Cref{subsec:mlp}), 1D~CNN
(\Cref{subsec:cnn1d}), or ResNet-1D (\Cref{subsec:resnet1d})
architectures described above, among many others.  Similarly,
scikit-learn~\cite{sklearn_docs} provides implementations of logistic
regression, SVM variants, Random Forest, Gradient Boosted Trees, MLP,
and K-Means.  The choice of implementation framework affects engineering
convenience, deployment compatibility, and edge-runtime availability,
but it does not change the mathematical properties of the underlying
algorithm.  Framework selection is discussed in \Cref{sec:nextsteps}
as part of the implementation stack.

% ==================================================================
\subsection{Level 1: Linear Baselines}\label{subsec:level1}
% ==================================================================
\subsubsection{Logistic Regression}\label{subsec:logreg}
% ==================================================================

\paragraph{Core mechanism.}
Multinomial logistic regression extends binary logistic regression to
the three-class setting by modelling the posterior probability of each
class as a softmax over linear combinations of the input
features~\cite{hosmer2013logistic}.  For an input vector
$\mathbf{x} \in \mathbb{R}^{d}$ and $K=3$ classes, the model computes
%
\begin{equation}
  P(y = k \mid \mathbf{x})
  = \frac{\exp(\mathbf{w}_k^\top \mathbf{x} + b_k)}
         {\sum_{j=1}^{K} \exp(\mathbf{w}_j^\top \mathbf{x} + b_j)},
  \qquad k = 1, \dots, K,
\end{equation}
%
where $\mathbf{w}_k$ and $b_k$ are the weight vector and bias for
class~$k$.  The decision boundaries are hyperplanes in the
$d$-dimensional feature space.

\paragraph{Expected input shape.}
A flattened one-dimensional feature vector of length $d$.  With the
three-channel hybrid encoding (dB, Re, Im), $d = 12003$; ablation
variants use $d = 4001$ (dB only) or $d = 8002$ (Re, Im only).

\paragraph{Compatible route.}
Route~A exclusively.  The model has no mechanism to exploit
two-dimensional spatial structure.

\paragraph{Required preprocessing.}
Standard shared pipeline from \Cref{sec:preproc}: Touchstone parsing,
frequency-grid validation, class-label assignment.

\paragraph{Preferred normalization.}
Per-feature z-score normalization (fit on training partition only) is
preferred because the magnitude of $S_{11}$ in dB and phase in degrees
occupy different numerical ranges.  Min--max scaling is acceptable but
more sensitive to outliers.

\paragraph{Augmentation compatibility.}
Compatible with all Route~A augmentations: additive Gaussian noise,
magnitude scaling, phase perturbation.  Since the model has minimal
capacity, augmentation primarily serves as a regulariser against
noise rather than a capacity enhancer.

\paragraph{Overfitting risk.}
\textbf{Very low.}  The total parameter count is $K \times (d + 1)$,
which is at most $3 \times 8003 = 24{,}009$ parameters --- well below
the $\num{1603}$-sample regime where high-capacity models become
dangerous.  $L_2$ regularisation (controlled by the inverse
regularisation strength $C$) further constrains the weight magnitudes.

\paragraph{Interpretability.}
\textbf{High.}  Each weight $w_{k,i}$ directly quantifies the
contribution of frequency bin~$i$ to the log-odds of class~$k$.  This
enables direct mapping of learned coefficients onto the physical
frequency axis, revealing which spectral regions drive classification.

\paragraph{Edge deployment suitability.}
\textbf{Excellent.}  The inference operation is a single matrix--vector
multiplication followed by a softmax, requiring $\mathcal{O}(Kd)$
multiply-accumulate operations.  The model fits in less than
\SI{100}{\kilo\byte} even without quantisation.  No specialised
runtime is required; inference can be implemented in plain~C.

\paragraph{Main failure modes.}
Logistic regression assumes that the classes are linearly separable (or
approximately so) in the feature space.  If the physical differences
between \texttt{wet} and \texttt{ice} manifest as nonlinear patterns
--- for instance, resonance peak \emph{shifts} combined with
magnitude \emph{dips} at different frequency locations --- a linear
boundary will misclassify the overlap region.  Dimensionality reduction
to a small number of principal components may also discard discriminative
spectral detail.


% ==================================================================
\subsubsection{Linear SVM}\label{subsec:lsvm}
% ==================================================================

\paragraph{Core mechanism.}
A support vector machine with a linear kernel finds the maximum-margin
separating hyperplane between classes~\cite{cortes1995svm,
vapnik1995svm}.  For three classes, the implementation typically uses
a one-vs-rest (OVR) or one-vs-one (OVO) decomposition, training
$K$ or $\binom{K}{2}$ binary classifiers respectively.

\paragraph{Expected input shape.}
Flattened 1D vector, same as logistic regression.

\paragraph{Compatible route.}
Route~A exclusively.

\paragraph{Required preprocessing.}
Shared pipeline from \Cref{sec:preproc}.

\paragraph{Preferred normalization.}
Z-score or min--max scaling.  The margin maximisation objective is
sensitive to feature scale; un-normalised features with large dynamic
range will dominate the margin calculation and suppress contributions
from lower-magnitude features.

\paragraph{Augmentation compatibility.}
Same as logistic regression.  Noise-based augmentations can be viewed
as implicit regularisation of the margin.

\paragraph{Overfitting risk.}
\textbf{Low.}  The margin constraint is itself a structural
regulariser.  The penalty parameter $C$ controls the trade-off between
margin width and training-set misclassification tolerance.  With proper
cross-validated selection of~$C$, the model generalises well even on
small datasets.

\paragraph{Interpretability.}
\textbf{High.}  The normal vector of each separating hyperplane can be
projected back onto the frequency axis, analogously to logistic
regression coefficients.  Support vectors identify the most ambiguous
training samples near the decision boundary.

\paragraph{Edge deployment suitability.}
\textbf{Good.}  Inference requires computing the sign of
$\mathbf{w}^\top \mathbf{x} + b$ for each binary classifier.  For
OVR with $K=3$, this is three inner products.  If the support vector
set is sparse (common for well-separated data), the stored model is
very compact.

\paragraph{Main failure modes.}
Identical to logistic regression: linear decision boundaries are
insufficient if class separations are inherently nonlinear.  The linear
SVM may achieve slightly different margins than logistic regression on
the same data, but neither can overcome the fundamental limitation of
linearity.



% ==================================================================
\subsection{Level 2: Non-Linear Tabular Models}\label{subsec:level2}
% ==================================================================
\subsubsection{RBF SVM}\label{subsec:rbfsvm}
% ==================================================================

\paragraph{Core mechanism.}
Replacing the linear kernel with the Radial Basis Function (RBF)
kernel
$K(\mathbf{x}, \mathbf{x}') = \exp\!\bigl(-\gamma \|\mathbf{x} -
\mathbf{x}'\|^2\bigr)$
implicitly maps the input to an infinite-dimensional Hilbert space,
enabling nonlinear decision boundaries in the original feature
space~\cite{vapnik1995svm}.

\paragraph{Expected input shape.}
Flattened 1D vector, same as the linear methods.

\paragraph{Compatible route.}
Route~A.  Although the kernel can, in principle, operate on any fixed-length
vector, its distance-based nature assumes that all feature dimensions are
commensurable after normalisation.

\paragraph{Required preprocessing.}
Shared pipeline from \Cref{sec:preproc}.

\paragraph{Preferred normalization.}
\textbf{Z-score normalization is essential.}  The RBF kernel computes
Euclidean distances between feature vectors; features on larger scales
will dominate these distances and distort the kernel matrix.  Robust
scaling is an alternative when extreme outliers are suspected.

\paragraph{Augmentation compatibility.}
Compatible with all Route~A augmentations.  Because the RBF SVM can
already model nonlinear boundaries, augmentation primarily improves
robustness to measurement noise rather than increasing effective
complexity.

\paragraph{Overfitting risk.}
\textbf{Medium.}  Two hyperparameters require careful tuning:
$\gamma$ (kernel width) and $C$ (penalty).  A high $\gamma$ combined
with a large $C$ can produce a decision surface that tightly wraps
around individual training points, effectively memorising the
training set.  Nested cross-validation is essential to control this
risk.

\paragraph{Interpretability.}
\textbf{Low--medium.}  The decision function depends on kernel
evaluations against the support vectors, making it difficult to trace
predictions back to individual frequency bins.  Partial dependence
plots or SHAP-based kernel explanations can provide post-hoc
interpretability, but at additional computational cost.

\paragraph{Edge deployment suitability.}
\textbf{Moderate.}  Inference requires computing
$f(\mathbf{x}) = \sum_{i \in \mathcal{SV}} \alpha_i y_i
K(\mathbf{x}, \mathbf{x}_i) + b$,
where the sum is over all support vectors.  If the number of support
vectors is large (a risk with noisy or overlapping classes), both
memory and computation scale unfavourably.  Approximate methods (e.g.,
Nyström approximation) or explicit kernel-map expansions can mitigate
this but add implementation complexity.

\paragraph{Main failure modes.}
Poor scaling to very high-dimensional inputs without prior
dimensionality reduction.  With $d = 4001$ or $d = 8002$, the pairwise
kernel matrix becomes expensive to compute during training, and the
support vector set may become dense.  Additionally, the kernel
hyperparameter $\gamma$ can be fragile: values that work well for one
noise level may fail under distribution shift.


% ==================================================================
\subsubsection{Random Forest}\label{subsec:rf}
% ==================================================================

\paragraph{Core mechanism.}
A Random Forest is an ensemble of $B$ decision trees, each trained on
a bootstrap sample of the training data with random feature subsampling
at each split~\cite{breiman2001rf}.  The ensemble prediction is the
majority vote (classification) or average (probability estimate) across
all trees.  Feature subsampling decorrelates the individual trees,
reducing variance relative to a single deep tree.

\paragraph{Expected input shape.}
Flattened 1D feature vector.  Each feature corresponds to a frequency
bin (magnitude, phase, or derived quantity).

\paragraph{Compatible route.}
Route~A.  While decision trees can technically operate on any tabular
input, they have no mechanism to exploit 2D spatial structure.

\paragraph{Required preprocessing.}
Shared pipeline from \Cref{sec:preproc}.

\paragraph{Preferred normalization.}
\textbf{Not strictly required.}  Decision-tree splits compare each
feature against a threshold; these comparisons are invariant to
monotonic transformations of the features.  However, if downstream
interpretation or ensemble combination with normalisation-sensitive
models is planned, applying z-score normalisation for consistency is
harmless.

\paragraph{Augmentation compatibility.}
Compatible with Route~A augmentations.  Noise injection can provide
additional regularisation, but the bootstrap mechanism already
introduces sampling diversity.

\paragraph{Overfitting risk.}
\textbf{Low--medium.}  Controlled by the number of trees
(\texttt{n\_estimators}), maximum tree depth (\texttt{max\_depth}),
and minimum samples per leaf (\texttt{min\_samples\_leaf}).
Unrestricted depth allows each tree to perfectly memorise its bootstrap
sample; setting \texttt{max\_depth} $\leq 20$ or
\texttt{min\_samples\_leaf} $\geq 5$ is recommended for $\num{1603}$
samples.  Increasing \texttt{n\_estimators} does not increase
overfitting risk (it reduces variance), but it increases model size.

\paragraph{Interpretability.}
\textbf{Good.}  Feature importance scores (mean decrease in impurity
or permutation importance) map directly to frequency bins, revealing
which spectral regions are most discriminative.  This provides a
physically interpretable summary of which portions of the $S_{11}$
trace carry classification information.

\paragraph{Edge deployment suitability.}
\textbf{Good.}  Each tree is a sequence of if--else comparisons with
no matrix operations.  Trees can be serialised as compact lookup
structures.  Inference can be parallelised across trees.  However,
large ensembles ($B > 100$) with deep trees may exceed the memory
budget of constrained microcontrollers.  Ensemble pruning (removing
low-contribution trees) is straightforward.

\paragraph{Main failure modes.}
Individual trees can memorise the training set if depth is
unrestricted, but the ensemble average typically smooths this out.
The primary practical risk is \emph{model size}: an ensemble of 500
deep trees with $d = 8002$ features can produce a serialised model
exceeding \SI{50}{\mega\byte}, which is prohibitive for edge
deployment.  Feature subsampling with $\sqrt{d} \approx 90$ features
per split keeps individual trees lightweight, but total ensemble
memory must be profiled.  Random Forests can also underperform when
the discriminative signal is spread thinly across many correlated
features rather than concentrated in a few strong predictors.


% ==================================================================
\subsubsection{Gradient Boosted Trees}\label{subsec:gbt}
% ==================================================================

\paragraph{Core mechanism.}
Gradient Boosted Trees (GBT) build an additive ensemble of shallow
decision trees in a sequential (boosting) fashion, where each new tree
is trained to correct the residual errors of the current
ensemble~\cite{friedman2001gbm}.  Modern implementations such as
XGBoost~\cite{chen2016xgboost} add column subsampling,
second-order gradient information, and regularisation terms to the
objective, making GBT one of the most competitive methods for
tabular data across many application domains.

\paragraph{Expected input shape.}
Flattened 1D feature vector, identical to Random Forest.

\paragraph{Compatible route.}
Route~A.

\paragraph{Required preprocessing.}
Shared pipeline from \Cref{sec:preproc}.

\paragraph{Preferred normalization.}
Not strictly required for the same reasons as Random Forest (tree
splits are scale-invariant), but normalisation is harmless and aids
pipeline consistency.

\paragraph{Augmentation compatibility.}
Compatible with Route~A augmentations.  Column subsampling and row
subsampling within XGBoost already provide built-in stochastic
regularisation; additional augmentation can complement but should
be applied judiciously to avoid excessive noise.

\paragraph{Overfitting risk.}
\textbf{Medium.}  The sequential nature of boosting means that each
additional tree further reduces training loss.  Without early stopping,
GBT can drive the training loss to zero and memorise the training set.
Three hyperparameters jointly control capacity: \texttt{learning\_rate}
(shrinkage), \texttt{max\_depth} (per-tree complexity), and
\texttt{n\_estimators} (total rounds).  A low learning rate
($\eta \leq 0.1$) with early stopping on a validation set is the
standard remedy.  Nested cross-validation is recommended to guard
against optimistic tuning on a single split.

\paragraph{Interpretability.}
\textbf{Good.}  Like Random Forest, GBT provides feature importance
scores that map to frequency bins.  SHAP (SHapley Additive
exPlanations) values are particularly well-developed for
tree-ensemble models and can provide both global and per-sample
explanations.

\paragraph{Edge deployment suitability.}
\textbf{Good.}  Boosted tree ensembles are typically shallower than
Random Forest trees (depth~3--6 is common), producing smaller models.
After pruning and quantisation, a trained GBT model can often fit
within \SI{1}{\mega\byte}.  Inference is a sequence of if--else
evaluations with no floating-point matrix operations, which is ideal
for microcontrollers without hardware FPU.

\paragraph{Main failure modes.}
Sensitivity to the learning rate: too high causes oscillation and
overfitting; too low requires many boosting rounds, inflating model
size and training time.  GBT can also overfit when the number of
informative features is small relative to the total feature count,
as uninformative features are repeatedly sampled and split upon.
Feature selection or dimensionality reduction before training can
mitigate this.



% ==================================================================
\subsection{Level 3: Deep Architectures}\label{subsec:level3}
% ==================================================================
\subsubsection{Multilayer Perceptron}\label{subsec:mlp}
% ==================================================================

\paragraph{Core mechanism.}
A Multilayer Perceptron (MLP) is a fully connected feedforward neural
network consisting of one or more hidden layers with nonlinear
activation functions~\cite{haykin2009nn, goodfellow2016deep}.  For an
input $\mathbf{x} \in \mathbb{R}^d$, a single-hidden-layer MLP
computes
%
\begin{equation}
  \hat{\mathbf{y}}
  = \mathrm{softmax}\!\bigl(\mathbf{W}_2 \,
    \sigma(\mathbf{W}_1 \mathbf{x} + \mathbf{b}_1) + \mathbf{b}_2\bigr),
\end{equation}
%
where $\sigma$ is a pointwise nonlinearity (e.g., ReLU), and
$(\mathbf{W}_1, \mathbf{b}_1, \mathbf{W}_2, \mathbf{b}_2)$ are
learnable parameters.  Deeper architectures stack additional
hidden layers.

\paragraph{Expected input shape.}
Flattened 1D vector of length $d = 12003$ (three-channel hybrid encoding,
\Cref{sec:repr:routeA:flat}) for Route~A inputs.  Ablation variants with
$d = 4001$ (dB only) or $d = 8002$ (Re, Im only) should be compared.
Route~B inputs (2D images) can also be flattened, although this discards
spatial structure.  Thus, the MLP is compatible with both routes in
principle but benefits from neither's structure.

\paragraph{Compatible route.}
Route~A (primary) or Route~B (via flattening).

\paragraph{Required preprocessing.}
Shared pipeline from \Cref{sec:preproc}.

\paragraph{Preferred normalization.}
Z-score or min--max normalisation is required.  Neural networks
trained with gradient descent are sensitive to feature scale:
un-normalised inputs cause poorly conditioned loss landscapes, slow
convergence, and potential numerical instability~\cite{goodfellow2016deep}.

\paragraph{Augmentation compatibility.}
Compatible with all Route~A augmentations and, if flattened Route~B
is used, with 2D augmentations applied before flattening.  Dropout
within the network provides implicit augmentation via stochastic
feature masking.

\paragraph{Overfitting risk.}
\textbf{Medium--high.}  A two-hidden-layer MLP with 256 and 128 units
has approximately $256d + 256 + 128 \times 256 + 128 + 3 \times 128 +
3 \approx 1{,}057{,}000$ parameters for $d = 4001$ --- roughly 660
parameters per training sample.  This ratio warrants aggressive
regularisation: dropout ($p = 0.3$--$0.5$), weight decay
($\lambda = 10^{-4}$--$10^{-3}$), and early stopping based on
validation loss.  Smaller architectures (e.g., 64--32 units) trade
capacity for safety.

\paragraph{Interpretability.}
\textbf{Low--medium.}  Unlike logistic regression, the hidden layers
introduce nonlinear transformations that obscure the relationship
between input features and predictions.  Gradient-based saliency maps
or integrated gradients can provide post-hoc attribution to input
frequency bins, but these are approximate and can be misleading
without careful validation.

\paragraph{Edge deployment suitability.}
\textbf{Good, if architecture is compact.}  A small MLP
(e.g., $4001 \to 64 \to 32 \to 3$) has approximately
\num{258000} parameters, requiring $\sim$\SI{1}{\mega\byte} in
float32 or $\sim$\SI{256}{\kilo\byte} in INT8.  Inference is
dominated by matrix--vector multiplications that are well-supported
by TensorFlow Lite and ONNX Runtime on ARM
targets~\cite{tensorflow_lite, onnx_runtime}.  Larger architectures
quickly exceed edge memory budgets.

\paragraph{Main failure modes.}
The MLP has \emph{no spectral or spatial inductive bias}: it treats
every input position independently, so there is no built-in notion
that adjacent frequency bins are related.  This means the MLP must
learn locality from scratch, which is data-inefficient.  With only
\num{1603} samples, the MLP may fail to discover localized spectral
patterns (e.g., resonance peak shifts) that a convolutional
architecture would capture natively.


% ==================================================================
\subsubsection{1D Convolutional Neural Network}\label{subsec:cnn1d}
% ==================================================================

\paragraph{Core mechanism.}
A 1D Convolutional Neural Network (CNN) applies learned convolutional
filters (kernels) along the frequency axis of the input
trace~\cite{lecun1998cnn}.  Each filter produces a feature map that
captures local patterns at its receptive field scale.  Stacking
convolutional layers with pooling operations builds a hierarchy of
increasingly abstract spectral features.

\paragraph{Expected input shape.}
A 1D sequence of length $L = 4001$ with $C = 3$ channels (dB, Re, Im;
\Cref{sec:repr:routeA:hybrid}).  Ablation variants with $C = 1$ (dB
only) or $C = 2$ (Re, Im only) should be included in the experimental
matrix.

\paragraph{Compatible route.}
Route~A (primary).  Multi-channel 1D input naturally accommodates
the three-channel representation discussed in \Cref{sec:repr}.

\paragraph{Required preprocessing.}
Shared pipeline from \Cref{sec:preproc}.

\paragraph{Preferred normalization.}
Z-score normalisation of the input, optionally supplemented by batch
normalisation layers within the network.  Batch normalisation
stabilises training by normalising activations at each layer,
reducing sensitivity to the input scale~\cite{goodfellow2016deep}.

\paragraph{Augmentation compatibility.}
Highly compatible with all Route~A augmentations.  The convolutional
architecture is particularly well suited to benefit from frequency-local
jitter and additive noise augmentations, as these diversify the
local patterns that the kernels must learn to recognise.

\paragraph{Overfitting risk.}
\textbf{Medium.}  Convolutional layers share parameters across
positions, drastically reducing the parameter count relative to an
MLP of comparable representational capacity.  A compact 1D CNN
(e.g., three convolutional blocks with 16--32--64 filters followed by
global average pooling and a dense head) may have fewer than
\num{50000} parameters.  However, overfitting remains a concern with
\num{1603} samples; dropout, weight decay, and early stopping are
recommended.

\paragraph{Spectral inductive bias.}
This is the key advantage of the 1D CNN for $S_{11}$ classification.
Convolutional kernels of width $k$ (e.g., $k = 5$--$11$) learn to
detect localized frequency-domain patterns: resonance peaks,
absorption dips, slope changes, and inflection points.  These are
precisely the features that distinguish \texttt{dry}, \texttt{wet},
and \texttt{ice} states in $S_{11}$ spectra.  The inductive bias
encodes the physically motivated assumption that adjacent frequency
bins carry correlated information.

\paragraph{Three-channel justification for CNN.}
The three-channel hybrid input ($C = 3$) is particularly well-suited
for 1D CNNs.  Channel~1 (dB) exposes the deep resonance dip
at~$\sim\SI{-62}{\decibel}$ as a high-contrast feature that standard
filters can detect without needing to learn a logarithmic transform.
Channels~2 and~3 (Re, Im) provide smooth, continuous signals that avoid
false edge detections caused by phase wrapping.  Together, the three
channels give the convolutional layers complementary views of the
same spectrum.

\paragraph{Depthwise and grouped convolutions (recommended ablation).}
Because the three channels have fundamentally different scales and
physical semantics (dB vs.\ dimensionless linear), an architecture
variant using \emph{depthwise separable} or \emph{grouped}
convolutions ($\text{groups} = C$) may outperform standard
convolutions that mix channels from the first layer.  This variant
processes each channel independently in the early layers and fuses
them via a pointwise ($1 \times 1$) convolution afterward.  This
should be evaluated as an ablation, not mandated as default.

\paragraph{Interpretability.}
\textbf{Medium.}  First-layer convolutional filters can be visualised
directly as frequency-domain templates.  Gradient-weighted class
activation maps (Grad-CAM, adapted to 1D) can highlight which spectral
regions most influence each class prediction, providing a physically
meaningful attribution.

\paragraph{Edge deployment suitability.}
\textbf{Good.}  1D convolutions are computationally efficient
($\mathcal{O}(L \times C_{\text{in}} \times C_{\text{out}} \times k)$
per layer) and are natively supported by TensorFlow Lite and ONNX
Runtime~\cite{tensorflow_lite, onnx_runtime}.  INT8 quantisation of
convolutional layers is well-studied and typically preserves accuracy
with minimal degradation.  A compact 1D CNN can achieve sub-millisecond
inference on ARM Cortex-A class processors.

\paragraph{Main failure modes.}
The receptive field of a single convolutional layer is limited to $k$
frequency bins.  Without sufficient depth, striding, or pooling, the
network may fail to capture \emph{global} spectral shape (e.g., the
overall tilt of $|S_{11}|$ across the full \SI{8}{}--\SI{12}{\giga\hertz}
band).  Architectures must be designed to ensure that the effective
receptive field covers the full frequency span before the classification
head.  Additionally, very deep 1D CNNs without residual connections may
suffer from vanishing gradients.


% ==================================================================
\subsubsection{ResNet-1D}\label{subsec:resnet1d}
% ==================================================================

\paragraph{Core mechanism.}
ResNet-1D adapts the residual network architecture~\cite{he2016resnet}
to one-dimensional inputs.  Each residual block computes
$\mathbf{h}_{\ell+1} = \mathbf{h}_\ell + \mathcal{F}(\mathbf{h}_\ell)$,
where $\mathcal{F}$ is a stack of convolutional, batch-normalisation,
and activation layers.  The skip (identity) connection enables gradient
flow through very deep networks, allowing the model to learn
hierarchical spectral features across multiple scales.

\paragraph{Expected input shape.}
1D sequence of length $L = 4001$ with $C = 3$ channels, identical to the
1D CNN.  Depthwise/grouped convolution ablations apply equally.

\paragraph{Compatible route.}
Route~A.

\paragraph{Required preprocessing.}
Shared pipeline from \Cref{sec:preproc}.

\paragraph{Preferred normalization.}
Z-score input normalisation combined with batch normalisation within
residual blocks.

\paragraph{Augmentation compatibility.}
Highly compatible with all Route~A augmentations.  Given the high
capacity of ResNet-1D, augmentation is not merely beneficial but
\emph{essential} for regularisation under the \num{1603}-sample
constraint.

\paragraph{Overfitting risk.}
\textbf{High.}  This is the highest-capacity model in the comparison.
A modest ResNet-1D with 8 residual blocks and 64 filters per block
can easily exceed \num{500000} parameters --- more than 300 parameters
per training sample.  Without aggressive regularisation, the model
will memorise the training set.  The following safeguards are mandatory:
%
\begin{itemize}
  \item \textbf{Dropout} ($p = 0.3$--$0.5$) after residual blocks.
  \item \textbf{Weight decay} ($\lambda = 10^{-3}$--$10^{-2}$).
  \item \textbf{Early stopping} with patience of 10--20 epochs on
        validation loss.
  \item \textbf{Architecture minimisation}: use the smallest depth
        and width that achieve competitive validation performance.
  \item \textbf{Data augmentation} at every training epoch.
\end{itemize}

\paragraph{Interpretability.}
\textbf{Low--medium.}  The depth of the architecture makes it harder
to trace predictions back to input features compared with a shallow
1D CNN.  Grad-CAM and saliency methods remain applicable but become
less reliable as network depth increases.

\paragraph{Edge deployment suitability.}
\textbf{Moderate.}  A well-pruned and quantised ResNet-1D can be
deployed on ARM Cortex-A class hardware via TensorFlow Lite or ONNX
Runtime~\cite{tensorflow_lite, onnx_runtime}, but may exceed the
constraints of Cortex-M class microcontrollers.  The residual
connections themselves add negligible overhead (element-wise addition),
so inference cost is dominated by the convolutional layers.  INT8
quantisation and channel pruning are viable reduction strategies.

\paragraph{Main failure modes.}
ResNet-1D is the model \emph{most prone to overfitting} in this
comparison.  The combination of high capacity and small training set
creates a setting where validation performance can diverge sharply
from training performance.  Careful architecture search (varying
depth, width, and regularisation simultaneously) is necessary, and
nested cross-validation is strongly recommended to obtain reliable
performance estimates.  Training instability (e.g., loss spikes) can
also occur if learning rates are not carefully scheduled.


% ==================================================================
\subsubsection{Kolmogorov--Arnold Network (KAN)}\label{subsec:kan}
% ==================================================================

\paragraph{Core mechanism.}
Kolmogorov--Arnold Networks~\cite{liu2024kan} replace the fixed
activation functions of traditional neural networks with
\emph{learnable univariate functions} placed on the edges (connections)
rather than the nodes.  This design is inspired by the
Kolmogorov--Arnold representation theorem~\cite{kolmogorov1957representation},
which states that any continuous multivariate function can be
decomposed into a finite composition of continuous univariate functions
and addition.  In practice, the learnable edge functions are
parameterised as B-spline curves with adjustable control points.

\paragraph{Expected input shape.}
Flattened 1D feature vector ($d = 12003$ for three-channel hybrid;
ablation variants with $d = 4001$ or $d = 8002$), Route~A.

\paragraph{Compatible route.}
Route~A.  Convolutional variants of KAN (KAN-Conv) have been proposed
but remain experimental.

\paragraph{Required preprocessing.}
Shared pipeline from \Cref{sec:preproc}.

\paragraph{Preferred normalization.}
Z-score normalisation is recommended, although the learnable activation
functions can, in principle, adapt to arbitrary input scales.  In
practice, well-scaled inputs accelerate convergence and improve
numerical stability of the B-spline parameterisation.

\paragraph{Augmentation compatibility.}
Compatible with Route~A augmentations, but the interaction between
augmented data and learned B-spline activations is not yet well
characterised in the literature.

\paragraph{Overfitting risk.}
\textbf{High.}  KAN introduces additional learnable parameters
(B-spline control points) on every edge, potentially exceeding the
parameter count of a comparable MLP.  The architecture is new, and
regularisation strategies tailored to KAN are still under active
research.  With \num{1603} samples, overfitting is a significant
concern.  Weight decay, grid-based sparsification, and early stopping
should be employed, but their effectiveness on small $S_{11}$ datasets
is unknown.

\paragraph{Interpretability.}
\textbf{Potentially high.}  A distinctive claimed advantage of KAN is
that the learned univariate functions on each edge can be individually
visualised and inspected.  If the network is sufficiently shallow
(e.g., two layers), each edge function represents a learned
transformation of a specific input feature or intermediate
representation.  This could, in principle, reveal physically meaningful
nonlinear relationships between specific frequency bins and the
classification outcome.  However, this interpretability has been
demonstrated primarily on low-dimensional synthetic tasks; its utility
on a $d = 4001$-dimensional spectral input is unvalidated.

\paragraph{Edge deployment suitability.}
\textbf{Uncertain.}  As of mid-2025, there is no established deployment
toolchain for KAN on embedded hardware.  B-spline evaluation requires
locating the correct interval and computing polynomial basis functions,
which is more complex per-parameter than the multiply-accumulate
operations of standard neural networks.  Neither TensorFlow Lite nor
ONNX Runtime natively support KAN layers.  Custom inference kernels
would need to be implemented and optimised, adding significant
engineering risk.

\paragraph{Main failure modes.}
The KAN ecosystem is immature.  Limited scaling studies exist for
inputs with $d > 100$.  Training can exhibit instability due to the
B-spline parameterisation (control-point oscillations, poor
initialisation).  The lack of established best practices for
architecture selection, regularisation, and hyperparameter tuning on
real-world datasets makes KAN a high-risk, high-potential-reward choice
in this comparison.  It is included for methodological completeness
and because its interpretability claims are particularly relevant to
physics-grounded sensing applications, but it should be treated as
experimental.


